% Authoritative RouterBSI_Draft_Appendix.docx; zero-based direct body paragraph IDs.

% DOCX-P000-BEGIN
\begin{abstract}
% DOCX-P000-END

% DOCX-P001-BEGIN
As agents are expected to solve broader and more complex tasks, they rely on external knowledge saved as skills to handle tasks beyond the LLM's parametric capabilities. As the skill corpora grows into hundreds or thousands of skills, skill routing becomes a bottleneck, where the agent must efficiently retrieve and select the appropriate skills to complete a task. Recent skill routing approaches leverage task feedback to improve skill routing performance online. However, feedback-driven updates to routing criteria can regress performance, and accumulating task feedback at scale requires efficient experience management. We introduce a recursive self-improving algorithm for skill routing in a black-box setting called RouterBSI. It refines task applicability predicates using Minimum Bayes Risk (MBR) with bounded replay to minimize errors on seen tasks and a held-out validation gate to reduce regressions on unseen tasks. The router evaluates predicates for top-k retrieved skills then selects skills if the task satisfies their respective predicate. RouterBSI improves its cold-start top-5 recall of required skills by 24.9\% across different retrievers on tasks from TheoremQA, LogicBench, ToolQA, MedCalc-Bench, CHAMP, and BigCodeBench sharing a 26k skill corpus. On average, it retains 90.5\% of the required skills in top 5 from the top 50 skills retrieved.
% DOCX-P001-END

\end{abstract}
